% =====================================================================
%  7. Monitoramento, avaliação e indicadores
% =====================================================================
\section{Monitoramento, avaliação e indicadores}
\label{sec:monitoramento}

O monitoramento acompanha a ação em dois níveis. No nível da execução, os
indicadores do \autoref{qua:indicadores} medem as entregas físicas: poços e
sistemas implantados, domicílios atendidos com água e com esgotamento e
população beneficiada. No nível dos resultados, a ação contribui para três
indicadores do Componente~3 (\autoref{tab:metas-componente3}): pessoas
atendidas com água gerida de forma segura, pessoas atendidas com pelo menos
esgotamento sanitário e associações de saneamento rural com modelo de gestão
adotado. Os valores de pessoas atendidas consideram 2,8 moradores por
domicílio \cite{ibge2022} e, no caso da água, incluem os sistemas implantados
pela Sanepar \cite{mop2026}.

\begin{quadro}[htbp]
\caption{Indicadores de implantação e reabilitação dos sistemas de abastecimento de água e esgotamento sanitário no meio rural}
\label{qua:indicadores}
\footnotesize
\renewcommand{\arraystretch}{1.25}
\begin{tabularx}{\textwidth}{|L|L|Q{1.6cm}|P{2.2cm}|Q{1.4cm}|}
\hline
\cab{Indicador} & \cab{Fórmula de cálculo} & \cab{Horizonte} &
\cab{Unidade} & \cab{Meta final} \\ \hline
Projetos de poços tubulares com sistemas de abastecimento de água implantados &
Quantidade de projetos executados & 1º ao 6º ano & Projetos (nº) & 83 \\ \hline
Domicílios atendidos com sistemas de abastecimento de água &
Soma das comunidades dentro dos municípios contemplados pelas contratações &
1º ao 6º ano & Domicílios & 5.810 \\ \hline
Domicílios atendidos com esgotamento sanitário &
70\% dos domicílios atendidos com abastecimento de água &
4º ao 6º ano & Domicílios & 4.067 \\ \hline
Domicílios atendidos com esgotamento sanitário nas propriedades rurais do IDR-Paraná &
Obtidos a partir dos dados do Censo 2022 & 4º ao 6º ano & Domicílios & 2.300 \\ \hline
População rural atendida com acesso à água potável e ao esgotamento sanitário &
Soma dos habitantes beneficiados pelas obras & 4º ao 6º ano &
Pessoas (nº) & 22.708 \\ \hline
\end{tabularx}
\fonte{Adaptado de \citeonline[Quadro 27]{mop2026}.}
\end{quadro}

\begin{table}[htbp]
\centering
\caption{Metas progressivas dos indicadores do Componente 3 relacionados ao Subcomponente 3.2.b}
\label{tab:metas-componente3}
\footnotesize
\renewcommand{\arraystretch}{1.2}
\begin{adjustbox}{max width=\linewidth}
\begin{tabular}{lrrrrrrr}
\toprule
\textbf{Indicador} & \textbf{2025} & \textbf{2027} & \textbf{2028} &
\textbf{2029} & \textbf{2030} & \textbf{2031} & \textbf{2032} \\
\midrule
Pessoas com água gerida de forma segura & 0 & 980 & 10.900 & 23.000 & 35.400 & 49.000 & 59.780 \\
\quad sexo feminino & 0 & 400 & 5.200 & 11.000 & 16.000 & 23.300 & 28.500 \\
\quad jovens (15 a 24 anos) & 0 & 110 & 1.300 & 2.800 & 4.200 & 5.900 & 7.000 \\
Pessoas com, pelo menos, esgotamento sanitário & 0 & 126 & 3.000 & 6.700 & 10.600 & 14.600 & 17.800 \\
\quad sexo feminino & 0 & 60 & 1.400 & 3.200 & 5.000 & 6.900 & 8.000 \\
\quad jovens (15 a 24 anos) & 0 & 15 & 360 & 800 & 1.200 & 1.700 & 2.000 \\
Associações RWSS com modelo de gestão adotado & 0 & 0 & 0 & 10 & 30 & 90 & 122 \\
\bottomrule
\end{tabular}
\end{adjustbox}
\fonte{Adaptado de \citeonline[Quadro 23]{mop2026}.}
\end{table}

\pendencia{Revisar as fórmulas do \autoref{qua:indicadores}. O indicador de
domicílios com abastecimento é calculado como ``soma das comunidades'', uma
unidade diferente de domicílios. O MOP escreve ``70\% do total de pessoas
atendidas domicílios atendidos com abastecimento'', o que mistura pessoas e
domicílios. O horizonte do esgotamento (4º ao 6º ano) não coincide com o
prazo das atividades (anos 2 a 5). O MOP também remete a um ``Quadro XX''.}

As informações que alimentam esses indicadores percorrem um fluxo definido
(\autoref{fig:monitoramento}). A supervisão e a fiscalização das obras
produzem, de forma contínua, relatórios de supervisão, laudos de vistoria e
testes de estanqueidade, validados pelo IAT. A cada semestre, o IAT
consolida essas informações, junto com o acompanhamento dos tempos de
resposta do Mecanismo de Queixas e Reclamações, em relatórios de
conformidade técnica e ambiental, cujos dados são registrados no SIGARH e no
GeoPR. Os indicadores de resultado têm frequência anual. O IAT e a Sanepar
reportam à Unidade de Gerenciamento do Programa (UGP) o número acumulado de
domicílios conectados e de associações que adotaram o modelo de gestão, a
partir de seus próprios registros e do banco de dados das associações. A UGP
consolida e valida essas informações à luz do MGAS e as reporta ao Banco
Mundial, que também valida, ao final, o relatório de encerramento e o Termo de
Recebimento Definitivo \cite{mop2026}.

\begin{figure}[htbp]
  \centering
  \caption{Fluxo de monitoramento, validação e reporte}
  \label{fig:monitoramento}
  \fluxograma{fluxograma-monitoramento}
  \fonte{Elaborado pelo IAT (2026) com base em \citeonline{mop2026}.}
\end{figure}
